\begin{theorem}[Precongruence for restriction]
  \label{thm:precong-relabel}
  \lean{PLTS.WeakScheduler.relabel}\lean{PLTS.ProbabilisticForwardSimulation.relabel}\leanok
  \uses{def:forwardsim, def:relabel, def:weakstep}
  If $sys_C$ is forward-simulated by $sys_A$ along $R$, both over an extended alphabet $\mathrm{Label} \oplus \mathrm{Extra}$, then $\mathrm{relabel}(sys_C)$ is forward-simulated by $\mathrm{relabel}(sys_A)$ along the same relation $R$.
\end{theorem}

\begin{proof}\leanok
  \uses{def:relabel, def:weakstep, def:weaktau}
  The initial condition is unchanged, restriction not touching initial states.
  A transition of $\mathrm{relabel}(sys_C)$ on $l$ is a transition of $sys_C$ on $\mathrm{inl}\,l$, which $R$ matches by a weak transition of $sys_A$; what has to be relabelled is the witness of that weak transition, a scheduler whose emissions are typed by the label type that has changed.
  Ask the witness about the prefix read on the extended alphabet and read its answer back along the map sending $\mathrm{inl}\,l$ to $l$ and every auxiliary label to $\tau$: the relabelled scheduler is valid, emits $\tau$ alone, and carries the original's mass on each emission and the original's halting mass, the auxiliary preimages of a base label all carrying mass $0$.
  The two weak transitions therefore have the same source and the same target distribution, the coupling is the one $R$ provides, and the case split between a silent and a visible matching transition carries over as it stands, $\mathrm{inl}\,\tau$ being the silent label of the extended alphabet.
\end{proof}
